% GENERATED FILE. Edit canonical lessons, then run npm run generate:books.
% canonical-source-hash: fnv1a64:11fee3f823897c41
% canonical-lessons: FR-C258-parlais, FR-C258-cetait, FR-C258-il-faisait-froid, FR-C258-dormais-dormi, FR-C258-pendant-que, FR-C258-revision-quand

\chapter{The Scene — c'était, il y avait, il faisait froid}
\label{ch:fr-the-scene}
\hlchaptermodality{258}

\begin{chapteropening}
\textbf{By the end of this chapter:} \emph{I can paint the scene of a past story with the imparfait, and drop the events into it with the passé composé.}

Complete FR-C258-revision-quand: a scene in the imparfait and an event that cuts into it, joined with quand or pendant que.
\end{chapteropening}

\section[je parlais, il parlait]{je parlais — the past that is still running}

\label{lesson:FR-C258-parlais}



\begin{quote}
Say \textquotedblleft{}I spoke\textquotedblright{}, with \emph{avoir} and the participle. (\emph{J'ai parlé.})

That past has an edge: it happened, and it finished. French has a second past with no edge at all, and every story you tell from here on needs both of them.
\end{quote}

\begin{grammarlens}[title={Grammar Lens: one stem, one set of endings}]
Start from the \emph{nous} form of the present and take off its \textbf{-ons}. What is left is the stem. Then add \textbf{-ais} for \textquotedblleft{}I\textquotedblright{}:

\begin{quote}
nous parl-ons, \textbf{je parlais} — I was speaking, I used to speak nous finiss-ons, \textbf{je finissais} — I was finishing nous pren-ons, \textbf{je prenais} — I was taking
\end{quote}

The other persons keep that stem and change only the ending: \textbf{je parlais}, \textbf{tu parlais}, \textbf{il parlait}, \textbf{nous parlions}, \textbf{vous parliez}, \textbf{ils parlaient}.

Two things to hear. Four of those six are one sound: \emph{parlais}, \emph{parlais}, \emph{parlait} and \emph{parlaient} all end in the same open \emph{è}, so the pronoun in front is what tells a listener who was speaking. And \emph{nous parlions} is one small \emph{i} away from the present \emph{nous parlons}. Listen for that \emph{i}.

This past is the \textbf{imparfait}, and almost every verb follows this recipe. The one that refuses it comes next.

What it means is \textquotedblleft{}was ‑ing\textquotedblright{}, or \textquotedblleft{}used to\textquotedblright{}: \emph{j'habitais à Paris}, I lived in Paris, back then, with no start and no finish in view.
\end{grammarlens}

\begin{cousinweb}
Latin already had this past, and it marked it with a \textbf{b}: \emph{habēbat}, \textquotedblleft{}he was having\textquotedblright{}, from \emph{habēre}, the Latin verb that gave \textbf{avoir} (EtymDB 2.1). In French the \emph{b} wore away entirely. You met one of its forms in the chapter on \emph{être}: \textbf{j'étais} goes back to \emph{stābam}, \textquotedblleft{}I was standing\textquotedblright{}.

Older French spelled the ending \textbf{-oit}, \emph{il parloit}, and the spelling \textbf{-ais} settled only in the nineteenth century, long after the sound had changed. Neither EtymDB nor Lewis \& Short records endings this late, so that history is the usual account in histories of French, given here as one.
\end{cousinweb}

\subsection*{Your turn}
Say these aloud:

\begin{itemize}
\raggedright
  \item \emph{nous parlons}, then \emph{je parlais}
  \item \emph{je parlais, tu parlais, il parlait} — one sound
  \item \emph{nous parlions, vous parliez} — and hear the \emph{i}
  \item \emph{j'habitais à Paris} — \emph{habiter}, still running in the past
  \item the edge, then no edge — \emph{j'ai parlé}, then \emph{je parlais}
\end{itemize}

\begin{tcolorbox}[breakable,colback=teal!4,colframe=teal!35!black,title={Before you move on}]
Which present form do you build the stem on? (The \emph{\textbf{nous}} form, without its \emph{-ons}.) Build \textquotedblleft{}I was taking\textquotedblright{}. (\emph{\textbf{Je prenais.}}) Which Latin letter wore away from the ending? (The \emph{\textbf{b}} of \emph{habēbat}.) Which past has an edge, and which has none? (\emph{\textbf{J'ai parlé}} has one; \emph{\textbf{je parlais}} has none.)
\end{tcolorbox}
\section[c'était, il y avait]{c'était, il y avait — the words a scene is made of}

\label{lesson:FR-C258-cetait}



\begin{quote}
Build \textquotedblleft{}I was speaking\textquotedblright{}. (\emph{Je parlais.}) Say \textquotedblleft{}it's good\textquotedblright{}. (\emph{C'est bon.}) Then \textquotedblleft{}there is bread\textquotedblright{}. (\emph{Il y a du pain.})

Hold on to all three. This lesson moves the last two into the past.
\end{quote}

\subsection*{You'll want to know: c'était}
One verb refuses last lesson's recipe: \textbf{être}. Its \emph{nous} form, \emph{nous sommes}, gives nothing to build on, so it has a stem of its own, \textbf{ét-}, with the same endings: \emph{j'étais}, \emph{tu étais}, \emph{il était}, \emph{nous étions}, \emph{vous étiez}, \emph{ils étaient}.

The form you will use most is the one inside \emph{c'est}:

\begin{quote}
\textbf{C'était} bon. — It was good. \textbf{C'était} samedi. — It was Saturday.
\end{quote}

\emph{C'est} is now; \textbf{c'était} is how things were, all through some stretch of the past.

\textbf{avoir} does follow the recipe: \emph{nous avons} gives \textbf{j'avais}. Age is the clearest case. \emph{J'ai vingt ans} is now; \emph{j'avais vingt ans} is the age you were, back then.

\subsection*{You'll want to know: il y avait}
You have \textbf{il y a}, \textquotedblleft{}there is\textquotedblright{} and \textquotedblleft{}there are\textquotedblright{}. Its past is \textbf{il y avait}:

\begin{quote}
\textbf{Il y avait} beaucoup de monde. — There were a lot of people. \textbf{Il n'y avait pas de} pain. — There was no bread.
\end{quote}

Like \emph{il y a}, it does not change for a plural: one \emph{avait} for one thing or for a hundred. And after the negative, \emph{pas de} stays exactly as it was in the present.

\begin{cousinweb}
\textbf{était} goes back to Latin \emph{stābat}, \textquotedblleft{}he was standing\textquotedblright{}, the standing verb that lent \emph{être} every form beginning \emph{ét-}. \textbf{avait} goes back to \emph{habēbat}, \textquotedblleft{}he was having\textquotedblright{}, the past of \emph{habēre}, the Latin verb that gave \textbf{avoir} (EtymDB 2.1). In both, the Latin \emph{b} wore away.

The little \textbf{y} of \emph{il y a} is inherited from Latin \emph{ibī}, \textquotedblleft{}there\textquotedblright{} (EtymDB 2.1; Lewis \& Short, under \emph{ibi}). So \emph{il y avait} is, word for word, \textquotedblleft{}it had there\textquotedblright{}.
\end{cousinweb}

\subsection*{Your turn}
Say these aloud:

\begin{itemize}
\raggedright
  \item \emph{c'était, il y avait}
  \item \emph{C'était samedi.} — then \emph{C'était bon.}
  \item \emph{J'avais vingt ans.}
  \item \emph{Il y avait beaucoup de monde.} — then \emph{Il n'y avait pas de pain.}
  \item from now to then — \emph{c'est} becomes \emph{c'était}, \emph{il y a} becomes \emph{il y avait}
\end{itemize}

\begin{tcolorbox}[breakable,colback=teal!4,colframe=teal!35!black,title={Before you move on}]
What stem does \emph{être} use in the imparfait? (\emph{\textbf{Ét-.}}) Say \textquotedblleft{}there were a lot of people\textquotedblright{}. (\emph{\textbf{Il y avait beaucoup de monde.}}) Does \emph{il y avait} change for a plural? (\textbf{No.}) Which Latin verb is behind \emph{était}? (\emph{\textbf{Stāre}}, \textquotedblleft{}to stand\textquotedblright{}.)
\end{tcolorbox}
\section[il faisait froid, il faisait chaud]{il faisait froid — the weather, made in the past}

\label{lesson:FR-C258-il-faisait-froid}



\begin{quote}
Say \textquotedblleft{}there were a lot of people\textquotedblright{}, the way last lesson built it. (\emph{Il y avait beaucoup de monde.}) Then \textquotedblleft{}it's cold\textquotedblright{}, today. (\emph{Il fait froid.})
\end{quote}

\subsection*{You'll want to know: il faisait froid, il faisait chaud}
English says the weather \textbf{was} cold. French says it \textbf{made} cold, and you have that already: \emph{il fait froid}. Move it into the past with the recipe. \emph{Nous faisons} gives the stem \emph{fais-}:

\begin{quote}
\textbf{Il faisait froid.} — It was cold. \textbf{Il faisait chaud.} — It was hot.
\end{quote}

Listen to the first syllable. The \emph{ai} of \emph{faisait} is said like the short \emph{e} of \emph{le}: \textquotedblleft{}fe-zè\textquotedblright{}. It is the sound of \emph{nous faisons} too, and the recipe carries it over.

Rain has a verb of its own. You have \textbf{il pleut}, \textquotedblleft{}it's raining\textquotedblright{}, and its past is \textbf{il pleuvait}, \textquotedblleft{}it was raining\textquotedblright{}.

Keep \textbf{c'était} and \textbf{était} for things. \emph{Le café était froid} — the coffee was cold — is about the coffee. \emph{Il faisait froid} is about the whole day. Put three of these together and you have a morning before anything has happened in it:

\begin{quote}
\emph{Il faisait froid et il pleuvait. Il n'y avait pas de soleil.}
\end{quote}

In Québec, relaxed speech has its own word for a biting cold: \emph{il faisait frette}. In a letter, \emph{froid} is the word to write.

\begin{cousinweb}
\textbf{faire} is inherited from Latin \emph{facere}, \textquotedblleft{}to make, to do\textquotedblright{} (EtymDB 2.1), and \textbf{faisait} goes back to its past \emph{faciēbat}, \textquotedblleft{}he was making\textquotedblright{} (Lewis \& Short, under \emph{facio}). English \textbf{fact} and \textbf{factory} share it.

\textbf{froid} is inherited, through Old French \emph{froit}, from Latin \emph{frīgidus}, \textquotedblleft{}cold\textquotedblright{} (EtymDB 2.1), which English borrowed whole as \textbf{frigid}. \textbf{chaud} goes back to Old French \emph{chaut} (EtymDB 2.1), usually traced to Latin \emph{calidus}, \textquotedblleft{}warm, hot\textquotedblright{} (Lewis \& Short).
\end{cousinweb}

\subsection*{Your turn}
Say these aloud:

\begin{itemize}
\raggedright
  \item \emph{Il faisait froid.} — then \emph{Il faisait chaud.}
  \item \emph{il pleut}, then \emph{il pleuvait}
  \item a thing, then a day — \emph{Le café était froid.} then \emph{Il faisait froid.}
  \item a whole morning — \emph{Il faisait froid et il pleuvait. Il n'y avait pas de soleil.}
\end{itemize}

\begin{tcolorbox}[breakable,colback=teal!4,colframe=teal!35!black,title={Before you move on}]
Say \textquotedblleft{}it was hot\textquotedblright{}. (\emph{\textbf{Il faisait chaud.}}) For the coffee, is it \emph{était} or \emph{faisait}? (\emph{\textbf{Était}} — it is about a thing.) How is the \emph{ai} of \emph{faisait} said? (Like the \emph{\textbf{e}} of \emph{le}.) Which Latin verb is inside \emph{faisait}? (\emph{\textbf{Facere}}, \textquotedblleft{}to make\textquotedblright{}.)
\end{tcolorbox}
\section[je dormais / j'ai dormi]{je dormais / j'ai dormi — the same sleep, told two ways}

\label{lesson:FR-C258-dormais-dormi}



\begin{quote}
Paint a cold morning in two short sentences. (\emph{Il faisait froid. Il n'y avait pas de soleil.})

Nothing has happened in it yet. This lesson is about what happens.
\end{quote}

\begin{grammarlens}[title={Grammar Lens: the scene and the events}]
Take one fact, that you slept, and let the tense decide what kind of fact it is:

\begin{quote}
\textbf{Je dormais.} — I was sleeping. \emph{(no edge: what was going on)} \textbf{J'ai dormi.} — I slept. \emph{(an edge: a thing that happened and ended)}
\end{quote}

Neither is more correct. A told story needs both, in a division of labour:

\begin{quote}
The \textbf{imparfait} paints the \textbf{scene}: what was going on, what things were like. The \textbf{passé composé} drops the \textbf{events} into it: what happened.
\end{quote}

\emph{Il faisait froid. Je dormais. Ma mère a téléphoné.} — It was cold. I was asleep. My mother phoned.

English has no separate ending for this. It leans on \emph{was ‑ing} and \emph{used to} and lets the context carry the rest. French settles it in the verb.
\end{grammarlens}

\begin{grammarlens}[title={Grammar Lens: what used to happen}]
A habit is scene too. It went on, again and again, with no single edge:

\begin{quote}
\emph{Autrefois, je dormais beaucoup.} — In the old days, I used to sleep a lot.
\end{quote}

So the imparfait answers two questions: what was going on at that moment, and what used to happen. The passé composé answers a third: what happened, once, and was done.
\end{grammarlens}

\begin{cousinweb}
\textbf{imparfait} is inherited from Latin \emph{imperfectus}, \textquotedblleft{}unfinished, incomplete\textquotedblright{} (EtymDB 2.1; Lewis \& Short): \emph{in-}, \textquotedblleft{}not\textquotedblright{}, plus \emph{perfectus}, \textquotedblleft{}carried all the way through\textquotedblright{}, which is itself \emph{per-}, \textquotedblleft{}through\textquotedblright{}, plus \emph{facere}, \textquotedblleft{}to make\textquotedblright{}. That \emph{facere} is the one inside last lesson's \emph{faisait}.

So the name describes the job. The imparfait is the \textbf{not-carried-through} past: unfinished in time, with no edge on it.
\end{cousinweb}

\subsection*{Your turn}
Say these aloud:

\begin{itemize}
\raggedright
  \item the pair — \emph{je dormais}, then \emph{j'ai dormi}
  \item a scene, then an event — \emph{Il pleuvait. Ma mère a téléphoné.}
  \item a habit — \emph{Autrefois, je dormais beaucoup.}
  \item \emph{C'était samedi, et j'ai dormi.}
\end{itemize}

\begin{tcolorbox}[breakable,colback=teal!4,colframe=teal!35!black,title={Before you move on}]
In a story, which past paints and which drops the events in? (The \emph{\textbf{imparfait}} paints; the \emph{\textbf{passé composé}} drops them in.) Is a habit scene or event? (\textbf{Scene}, in the imparfait.) What does \emph{imperfectus} say? (\textbf{Not carried through.})
\end{tcolorbox}
\section[quand, pendant que]{quand, pendant que — the scene, and what cut into it}

\label{lesson:FR-C258-pendant-que}



\begin{quote}
Say \textquotedblleft{}I was sleeping\textquotedblright{}, then \textquotedblleft{}I slept\textquotedblright{}. (\emph{Je dormais}, then \emph{j'ai dormi.})

This lesson joins a scene and an event into one sentence.
\end{quote}

\begin{grammarlens}[title={Grammar Lens: quand, where the event comes in}]
You have \textbf{quand}, \textquotedblleft{}when\textquotedblright{}. It puts a scene and an event into one sentence:

\begin{quote}
\emph{Je lisais \textbf{quand} ma mère a téléphoné.} — I was reading when my mother phoned. \emph{J'attendais le train \textbf{quand} il a commencé à pleuvoir.} — I was waiting for the train when it started to rain.
\end{quote}

The first half is the scene, still running, in the imparfait. The \emph{quand} half is the event that cut into it, in the passé composé. Turn the sentence round and the tenses keep their jobs: \emph{Quand ma mère a téléphoné, je lisais.}

A \emph{quand} can hold a scene too: \emph{Quand j'avais vingt ans, j'habitais à Paris.} A stretch of life is not an event, so both halves are scene.
\end{grammarlens}

\begin{grammarlens}[title={Grammar Lens: pendant que, two things at once}]
\textbf{pendant que} means \textquotedblleft{}while\textquotedblright{}. It holds a scene open so that something else can happen inside it:

\begin{quote}
\emph{\textbf{Pendant que} je lisais, mon frère dormait.} — While I was reading, my brother was asleep. \emph{\textbf{Pendant que} j'attendais le train, j'ai téléphoné à ma mère.} — While I was waiting for the train, I phoned my mother.
\end{quote}

Two scenes side by side, then one event inside a scene.

French has one more way to hold a moment still: \emph{\textbf{être en train de}}, \textquotedblleft{}to be in the middle of\textquotedblright{}. \emph{J'étais en train de lire quand ma mère a téléphoné.}
\end{grammarlens}

\begin{cousinweb}
\textbf{pendant} comes from \textbf{pendre}, \textquotedblleft{}to hang\textquotedblright{}, which is inherited from Latin \emph{pendēre}, \textquotedblleft{}to hang, to be suspended\textquotedblright{} (EtymDB 2.1; Lewis \& Short, under \emph{pendeo}). \emph{Pendant que je lisais} is, at root, \textquotedblleft{}while my reading was still hanging\textquotedblright{}: open, not yet settled. English \textbf{pending} has the same picture.

\textbf{quand} is inherited from Latin \emph{quandō}, \textquotedblleft{}when\textquotedblright{} (EtymDB 2.1), and the \emph{d} it kept in spelling is the Latin one.
\end{cousinweb}

\subsection*{Your turn}
Say these aloud:

\begin{itemize}
\raggedright
  \item \emph{Je lisais quand ma mère a téléphoné.}
  \item the same, turned round — \emph{Quand ma mère a téléphoné, je lisais.}
  \item \emph{Pendant que je lisais, mon frère dormait.}
  \item a scene of your own — \emph{Il faisait froid et j'attendais …} — then one thing that happened, with \emph{quand}
\end{itemize}

\begin{tcolorbox}[breakable,colback=teal!4,colframe=teal!35!black,title={Before you move on}]
In \emph{je lisais quand ma mère a téléphoné}, which half is the scene? (\emph{\textbf{Je lisais.}}) Which past does the \emph{quand} half take here? (The \textbf{passé composé} — the event.) What does \emph{pendant que} mean? (\textbf{While.})
\end{tcolorbox}
\section[(review)]{Review — the scene, and what cut into it}

\label{lesson:FR-C258-revision-quand}



\begin{quote}
Build \textquotedblleft{}I was reading\textquotedblright{} from \emph{lire}. (\emph{Je lisais.}) Then say \textquotedblleft{}I had a dog\textquotedblright{}. (\emph{J'avais un chien.})
\end{quote}

\begin{grammarlens}[title={What you've built}]
Five lessons, and every piece of a scene is in your hands. The recipe, the \emph{nous} stem plus \textbf{-ais}, on almost any verb: \emph{je parlais}, \emph{je finissais}, \emph{je prenais}. The one verb that refuses it, \emph{être}, with \emph{c'était}. The scene-setter \emph{il y avait}. The weather, made rather than being: \emph{il faisait froid}. And the division of labour: the imparfait paints, the passé composé drops the events in.

The joints between the two are \textbf{quand} and \textbf{pendant que}.

\begin{quote}
\emph{Je lisais quand ma mère a téléphoné.} \emph{J'attendais le train quand il a commencé à pleuvoir.}
\end{quote}

The second one is worth hearing twice. The rain is an event here, because it \textbf{began}: \emph{il a commencé}. Once it is falling, it is scene again: \emph{il pleuvait}.
\end{grammarlens}

\subsection*{Your turn}
Say these aloud:

\begin{itemize}
\raggedright
  \item a scene in three — \emph{Il faisait froid. Il y avait beaucoup de monde. C'était samedi.}
  \item \emph{Je lisais quand ma mère a téléphoné.}
  \item the rain as an event — \emph{il a commencé à pleuvoir} — then as a scene — \emph{il pleuvait}
  \item \emph{Pendant que j'attendais le train, j'ai téléphoné à ma mère.}
  \item one of your own — a scene with \emph{je …-ais}, and one thing that happened, with \emph{quand}
\end{itemize}

\begin{tcolorbox}[breakable,colback=teal!4,colframe=teal!35!black,title={Before you move on}]
After \emph{quand}, in these sentences, which past? (The \textbf{passé composé} — the event.) Why is \emph{il a commencé à pleuvoir} in the passé composé? (The rain \textbf{began}: that is an event.) What does \emph{imparfait} mean, taken apart? (\textbf{Not carried through.}) And which stem does \emph{c'était} use? (\emph{\textbf{Ét-}}, from the standing verb.)
\end{tcolorbox}
